% !TEX TS-program = lualatex
% !TEX encoding = UTF-8

\documentclass[VesperaleOP.tex]{subfiles}

\ifcsname preamble@file\endcsname
  \setcounter{page}{\getpagerefnumber{M-vop10_tempus_pentecosten}}
\fi

\begin{document}

%%%% IN DIE PENTECOSTES
\festum{}{In die Pentecostes}{I classis cum octava I classis}{In die Pentecostes}{1}

\titulumrubrica{Ad I Vesperas}
\cantus{P6F7VA}{Super Ps.}
\psalmus{145}
\psalmus{146}
\psalmus{147}

\capitulum{Act. 2}{Factus est repénte de cælo sonus tamquam adveniéntis spíritus veheméntis,\psstar{} et replévit totam domum ubi erant Apóstoli sedéntes.\pscross{}}
\cantus{P6F7VR}{\rrrub}
\cantus{P6F7VH}{Hymnus}

\versiculus{Repléti sunt omnes Spíritu Sancto, allelúia.}{Et cœpérunt loqui, allelúia.}

\cantus{P6F7VAM}{Ad Magn.}
\oratio{Præsta, quǽsumus, omnípotens Deus:\psstar{} ut claritátis tuæ super nos splendor effúlgeat;\pscross{}et lux tuæ lucis corda eórum, qui per grátiam tuam renáti sunt, Sancti Spíritus illustratióne confírmet.\psstar{} Per Dóminum... in unitáte eiúsdem.}

\titulumrubrica{Ad II Vesperas}
\cantus{P7F1VA}{Super Ps.}
\psalmus[3a]{109}
\psalmus[3a]{110}
\psalmus[3a]{111}

\rubrica{Capitulum, hymnus et \vvrub ut in I Vesperis.}

\cantus{P7F1VAM}{Ad Magn.}
\oratio{Deus, qui hodiérna die corda fidélium Sancti Spíritus illustratióne docuísti:\psstar{} da nobis in eódem Spíritu recta sápere,\pscross{}et de eius semper consolatióne gaudére.\psstar{} Per Dóminum... in unitáte eiúsdem.}
\lineamagna

\end{document}